% GENERATED FILE. Edit canonical lessons, then run npm run generate:books.
% canonical-source-hash: fnv1a64:9edef8bfa93ad7e4
% canonical-lessons: RU-C482-visok, RU-C482-zatylok, RU-C482-veko, RU-C482-nozdrya, RU-C482-chelyust

\chapter{Temple, Eyelid and Jaw}
\label{ch:ru-temple-eyelid-and-jaw}
\hlchaptermodality{482}

\begin{chapteropening}
\textbf{By the end of this chapter:} \emph{I can name the smaller parts of the head and face.}

Complete the last lesson of chapter 482: I can name the smaller parts of the head and face.
\end{chapteropening}

\section[\texorpdfstring{\ru{висок} (visók)}{висок (visók)}]{\ru{висок} (visók) — a temple (of the head)}

\label{lesson:RU-C482-visok}



\begin{quote}
Before the new one: say the Russian for a heel, then the Russian for a palm.
\end{quote}

\subsection*{You'll want to know: \ru{висок}}
\textbf{\ru{висок}} (\emph{visók}) — \textquotedblleft{}a temple (of the head)\textquotedblright{}.

\begin{quote}
\ru{У} \ru{него} \ru{седые} \ru{волосы} \ru{на} \ru{висках}.
\end{quote}

\begin{quote}
\emph{U negó sedýye vólosy na viskákh.}
\end{quote}

\begin{quote}
He has grey hair at the temples.
\end{quote}

Masculine; genitive \ru{виска}, plural \ru{виски} (stress on the end). Not to be confused with \emph{\ru{виски}} — whisky.

\begin{cousinweb}
EtymDB lists \emph{\ru{висок}} without a source; it is usually linked to \emph{\ru{висеть}} 'to hang', after the hair that hangs at the temples, but the link is not settled.
\end{cousinweb}

\begin{grammarlens}[title={What you've built}]
One more word for telling, explaining and describing at B1.
\end{grammarlens}

\subsection*{Your turn}
\begin{itemize}
\raggedright
  \item \emph{Say it:} \emph{visók}
  \item \emph{Say it:} \emph{visók}, once more
  \item \emph{Say it:} a sentence of your own with \emph{visók}
  \item \emph{Recall:} say the Russian for a heel, then the Russian for a palm, then say \emph{visók} again
\end{itemize}

\begin{tcolorbox}[breakable,colback=teal!4,colframe=teal!35!black,title={Before you move on}]
What does \ru{висок} mean? (\textquotedblleft{}A temple (of the head)\textquotedblright{}.) Say it once more.
\end{tcolorbox}
\section[\texorpdfstring{\ru{затылок} (zatýlok)}{затылок (zatýlok)}]{\ru{затылок} (zatýlok) — the back of the head}

\label{lesson:RU-C482-zatylok}



\begin{quote}
Before the new one: say the Russian for a palm, then the Russian for a temple.
\end{quote}

\subsection*{You'll want to know: \ru{затылок}}
\textbf{\ru{затылок}} (\emph{zatýlok}) — \textquotedblleft{}the back of the head\textquotedblright{}.

\begin{quote}
\ru{Мяч} \ru{попал} \ru{ему} \ru{в} \ru{затылок}.
\end{quote}

\begin{quote}
\emph{Myach popál yemú v zatýlok.}
\end{quote}

\begin{quote}
The ball hit him on the back of the head.
\end{quote}

Masculine; genitive \ru{затылка}.

\begin{cousinweb}
\emph{\ru{затылок}} is built on \emph{\ru{за}} 'behind' and \emph{\ru{тыл}} 'back, rear', which is inherited from Proto-Slavic \emph{tyl\ru{ъ}} 'back of the head, back' (EtymDB).
\end{cousinweb}

\begin{grammarlens}[title={What you've built}]
One more word for telling, explaining and describing at B1.
\end{grammarlens}

\subsection*{Your turn}
\begin{itemize}
\raggedright
  \item \emph{Say it:} \emph{zatýlok}
  \item \emph{Say it:} \emph{zatýlok}, once more
  \item \emph{Say it:} a sentence of your own with \emph{zatýlok}
  \item \emph{Recall:} say the Russian for a palm, then the Russian for a temple, then say \emph{zatýlok} again
\end{itemize}

\begin{tcolorbox}[breakable,colback=teal!4,colframe=teal!35!black,title={Before you move on}]
What does \ru{затылок} mean? (\textquotedblleft{}The back of the head\textquotedblright{}.) Say it once more.
\end{tcolorbox}
\section[\texorpdfstring{\ru{веко} (véko)}{веко (véko)}]{\ru{веко} (véko) — an eyelid}

\label{lesson:RU-C482-veko}



\begin{quote}
Before the new one: say the Russian for a temple, then the Russian for the back of the head.
\end{quote}

\subsection*{You'll want to know: \ru{веко}}
\textbf{\ru{веко}} (\emph{véko}) — \textquotedblleft{}an eyelid\textquotedblright{}.

\begin{quote}
\ru{У} \ru{меня} \ru{с} \ru{утра} \ru{дёргается} \ru{веко}.
\end{quote}

\begin{quote}
\emph{U menyá s utrá dyórgayetsya véko.}
\end{quote}

\begin{quote}
My eyelid has been twitching since the morning.
\end{quote}

Neuter; plural \ru{веки}.

\begin{cousinweb}
\emph{\ru{веко}} is inherited from Proto-Slavic \emph{věko} 'eyelid' (EtymDB).
\end{cousinweb}

\begin{grammarlens}[title={What you've built}]
One more word for telling, explaining and describing at B1.
\end{grammarlens}

\subsection*{Your turn}
\begin{itemize}
\raggedright
  \item \emph{Say it:} \emph{véko}
  \item \emph{Say it:} \emph{véko}, once more
  \item \emph{Say it:} a sentence of your own with \emph{véko}
  \item \emph{Recall:} say the Russian for a temple, then the Russian for the back of the head, then say \emph{véko} again
\end{itemize}

\begin{tcolorbox}[breakable,colback=teal!4,colframe=teal!35!black,title={Before you move on}]
What does \ru{веко} mean? (\textquotedblleft{}An eyelid\textquotedblright{}.) Say it once more.
\end{tcolorbox}
\section[\texorpdfstring{\ru{ноздря} (nozdryá)}{ноздря (nozdryá)}]{\ru{ноздря} (nozdryá) — a nostril}

\label{lesson:RU-C482-nozdrya}



\begin{quote}
Before the new one: say the Russian for the back of the head, then the Russian for an eyelid.
\end{quote}

\subsection*{You'll want to know: \ru{ноздря}}
\textbf{\ru{ноздря}} (\emph{nozdryá}) — \textquotedblleft{}a nostril\textquotedblright{}.

\begin{quote}
\ru{У} \ru{меня} \ru{заложена} \ru{правая} \ru{ноздря}.
\end{quote}

\begin{quote}
\emph{U menyá zalózhena právaya nozdryá.}
\end{quote}

\begin{quote}
My right nostril is blocked.
\end{quote}

Feminine; plural \ru{ноздри}.

\begin{cousinweb}
\emph{\ru{ноздря}} is inherited from Proto-Slavic \emph{nozdra}, which EtymDB traces to PIE \emph{néh\textsubscript{2}s 'nose', the root of }\ru{нос}\emph{.}
\end{cousinweb}

\begin{grammarlens}[title={What you've built}]
One more word for telling, explaining and describing at B1.
\end{grammarlens}

\subsection*{Your turn}
\begin{itemize}
\raggedright
  \item \emph{Say it:} \emph{nozdryá}
  \item \emph{Say it:} \emph{nozdryá}, once more
  \item \emph{Say it:} a sentence of your own with \emph{nozdryá}
  \item \emph{Recall:} say the Russian for the back of the head, then the Russian for an eyelid, then say \emph{nozdryá} again
\end{itemize}

\begin{tcolorbox}[breakable,colback=teal!4,colframe=teal!35!black,title={Before you move on}]
What does \ru{ноздря} mean? (\textquotedblleft{}A nostril\textquotedblright{}.) Say it once more.
\end{tcolorbox}
\section[\texorpdfstring{\ru{челюсть} (chélyust)}{челюсть (chélyust)}]{\ru{челюсть} (chélyust) — a jaw}

\label{lesson:RU-C482-chelyust}



\begin{quote}
Before the new one: say the Russian for an eyelid, then the Russian for a nostril.
\end{quote}

\subsection*{You'll want to know: \ru{челюсть}}
\textbf{\ru{челюсть}} (\emph{chélyust}) — \textquotedblleft{}a jaw\textquotedblright{}.

\begin{quote}
\ru{После} \ru{удара} \ru{у} \ru{него} \ru{болела} \ru{челюсть}.
\end{quote}

\begin{quote}
\emph{Pósle udára u negó boléla chélyust.}
\end{quote}

\begin{quote}
After the blow his jaw hurt.
\end{quote}

Feminine.

\begin{cousinweb}
\emph{\ru{челюсть}} is inherited from Proto-Slavic \emph{čeľust\ru{ь}} (EtymDB).
\end{cousinweb}

\begin{grammarlens}[title={What you've built}]
One more word for telling, explaining and describing at B1.
\end{grammarlens}

\subsection*{Your turn}
\begin{itemize}
\raggedright
  \item \emph{Say it:} \emph{chélyust}
  \item \emph{Say it:} \emph{chélyust}, once more
  \item \emph{Say it:} a sentence of your own with \emph{chélyust}
  \item \emph{Recall:} say the Russian for an eyelid, then the Russian for a nostril, then say \emph{chélyust} again
\end{itemize}

\begin{tcolorbox}[breakable,colback=teal!4,colframe=teal!35!black,title={Before you move on}]
What does \ru{челюсть} mean? (\textquotedblleft{}A jaw\textquotedblright{}.) Say it once more.
\end{tcolorbox}
